\problemstart{AP-02-010}{Compensated Summation Recovers Lost Increments}{Difficulty 3/5}{Chapter 2 \textbullet\ numpy}{Implement Kahan accumulation and measure order sensitivity against high-accuracy and library reductions.}

        \subsection{Mathematical statement}


Ordinary sequential accumulation applies
\[
  s_{k+1}=\operatorname{fl}(s_k+x_k).
\]
Kahan summation carries a compensation $c_k$ for low-order information:
\[
\begin{aligned}
y_k&=x_k-c_k,\\
t_k&=\operatorname{fl}(s_k+y_k),\\
c_{k+1}&=\operatorname{fl}((t_k-s_k)-y_k),\\
s_{k+1}&=t_k.
\end{aligned}
\]


        \subsection{Code problem}


Implement Kahan summation for a nonempty one-dimensional finite float64 array,
returning a Python float.  This explicit scalar loop is intentional: the problem
studies accumulation state.  Compare against Python \code{sum}, NumPy
\code{sum}, and \code{math.fsum}; note that since Python 3.12 the built-in
\code{sum} is itself Neumaier-compensated, so the naive baseline must be an
explicit $s_{k+1}=\operatorname{fl}(s_k+x_k)$ loop.  Do not claim Kahan is
correctly rounded for every adversarial sequence.


        \begin{contractbox}
        \begin{Code}
        def kahan_sum(values: np.ndarray) -> float:
        \end{Code}
        \end{contractbox}

        \subsection{Theory--engineering gap inventory}

        \begin{gapbox}
        \textbf{Explicit gap.} Associativity of real addition fails after rounding, so a reduction depends on order and algorithm despite denoting the same finite sum.

        \textbf{Implicit gap exposed by the result.} Many increments smaller than the current ULP vanish under sequential addition but reappear through compensation; reversed order changes the naive result.
        \end{gapbox}

        \subsection{Acceptance evidence}

        \begin{evidencebox}
        \begin{itemize}
          \item A small-increment fixture is closer to \code{math.fsum} than sequential Python accumulation.
  \item Simple exactly representable sums still match exactly.
  \item The benchmark prints signed errors and order sensitivity for each reduction separately. (Trust-based: the test suite does not machine-check benchmark output.)
        \end{itemize}
        \end{evidencebox}

        \subsection{Constraints}

        \begin{itemize}
          \item Use the stated Kahan recurrence; do not call \code{math.fsum} or \code{np.sum} inside the learner function.
  \item The scalar loop is pedagogically justified and needs no C extension at this scale.
  \item Reject empty, non-vector, non-float64, or nonfinite input with \code{ValueError}.
        \end{itemize}

        \begin{sourcebox}
        This is an atomic adaptation, not copied solution code. Nearest primary sources:
        \begin{itemize}
          \item \href{https://docs.python.org/3/library/math.html#math.fsum}{Python: math.fsum}
  \item \href{https://numpy.org/doc/stable/reference/generated/numpy.finfo.html}{NumPy: numpy.finfo}
        \end{itemize}
        \end{sourcebox}
